% Autorska priprava k Realisticky.cz 020312; ulohy jsou autorske a obmenene.
% Zamerne nenahrazuje vyklad presnymi kopiemi cizich textu/prikladu.
\documentclass[11pt,a4paper]{article}
\input{../../sablony/pripravy/priprava.sty}
\usepackage{tikz}
\nastavPripravu{1. ročník}{Rovnice a nerovnice}{83}{Soustavy tří rovnic o dvou neznámých}
\begin{document}
\blok{Návaznost a motivace}
Zatím jsme měli dvě neznámé a dvě rovnice. Co se stane, když přidáme třetí podmínku? Nemusí vzniknout „více řešení“; třetí rovnice může dosavadní řešení potvrdit, nebo je vyloučit.
\blok{Princip a odvození}
Nejdřív vyřešíme dvojici vhodných rovnic a nalezenou uspořádanou dvojici ověříme ve všech zbývajících rovnicích. Jsou-li dvě původní rovnice totožné, může být řešení stále nekonečně mnoho. Rozlišujeme počet rovnic a počet nezávislých podmínek.
\blok{Řešené příklady: od jednoduchého ke složitějšímu}
\begin{prikladReseniBox}[title={Příklad 1}]
Vyřešte $\begin{cases}x+y=5,\\2x-y=4,\\3x=9.\end{cases}$
\tcblower
\begin{align*}
3x&=9\Rightarrow x=3 \\
 y&=5-3=2 \\
 2x-y&=6-2=4
\end{align*}

Trojice rovnic má jediné společné řešení $K=\{[3;2]\}$. Třetí rovnice jej potvrzuje.
\end{prikladReseniBox}
\begin{prikladReseniBox}[title={Příklad 2}]
Vyřešte $\begin{cases}x+y=5,\\x-y=1,\\x+y=6.\end{cases}$
\tcblower
\begin{align*}
2x&=6\Rightarrow x=3 \\
 y&=2 \\
 x+y&=5\ne6
\end{align*}

Dvojice $[3;2]$ vyřeší první dvě rovnice, ale třetí odmítne. Proto $K=\varnothing$.
\end{prikladReseniBox}
\poznamka{Třetí rovnici nezanedbávej ani při zdánlivě hotovém výpočtu. Zdůrazni rozdíl mezi nezávislými a opakovanými podmínkami.}
\blok{Samostatné procvičování}
\begin{zadaniReseniBox}[title={Úloha 1}]
$\begin{cases}x+y=7,\\x-y=1,\\2x=8.\end{cases}$
\tcblower
\begin{align*}
2x&=8\Rightarrow x=4 \\
 y&=3 \\
 x-y&=1
\end{align*}

Všechny tři podmínky platí: $K=\{[4;3]\}$.
\end{zadaniReseniBox}
\begin{zadaniReseniBox}[title={Úloha 2}]
$\begin{cases}x+y=4,\\2x+2y=8,\\y=1.\end{cases}$
\tcblower
\begin{align*}
y&=1 \\
 x+1&=4\Rightarrow x=3 \\
 2x+2y&=6+2=8
\end{align*}
\end{zadaniReseniBox}
\begin{zadaniReseniBox}[title={Úloha 3}]
$\begin{cases}x+y=4,\\2x+2y=10,\\x-y=1.\end{cases}$
\tcblower
\begin{align*}
2(x+y)&=8\ne10 \\
 8&=10\;\text{-- spor} \\
 K&=\varnothing
\end{align*}
\end{zadaniReseniBox}
\blok{Volitelná rezerva}
\begin{zadaniReseniBox}[title={Úloha 1}]
Může mít soustava $x+y=5$, $x-y=1$ a $x=4$ společné řešení?
\tcblower
\begin{align*}
x+y&=5,\quad x-y=1\Rightarrow x=3,\ y=2 \\
 x&=4\;\text{-- nesplněno} \\
 K&=\varnothing
\end{align*}
\end{zadaniReseniBox}
\blok{Domácí úkol}
\begin{zadaniReseniBox}[title={Úloha 1}]
$\begin{cases}x+y=6,\\x-y=2,\\3x=12.\end{cases}$
\tcblower
\begin{align*}
2x&=8\Rightarrow x=4 \\
 y&=2 \\
 3x&=12
\end{align*}
\end{zadaniReseniBox}
\begin{zadaniReseniBox}[title={Úloha 2}]
$\begin{cases}2x+y=11,\\x-y=1,\\x+y=7.\end{cases}$
\tcblower
\begin{align*}
2x&=8\Rightarrow x=4 \\
 y&=3 \\
 2x+y&=11
\end{align*}
\end{zadaniReseniBox}
\begin{zadaniReseniBox}[title={Úloha 3}]
$\begin{cases}x+y=5,\\2x+2y=10,\\3x+3y=15.\end{cases}$
\tcblower
\begin{align*}
2(x+y)&=10 \\
 3(x+y)&=15 \\
 K&=\{[t;5-t]\mid t\in\mathbb R\}
\end{align*}
\end{zadaniReseniBox}
\poznamka{Orientační časové rozvržení: 5 + 12 + 12 + 10 + 6 = 45 min. Přednost má pochopení jádra; rezerva není součástí časového součtu. Domácí úlohy jsou určeny na samostatnou práci, řešení zůstávají učiteli.}
\end{document}
